\begin{textsample}{POS Dim 8 – global\_north\_2024\_07 – Score 11.00 – global\_north\_2024\_07/111183590.txt}  \label{ex:f8_pos_047}
\textbf{Vice} President Kamala Harris ' s decision to skip Prime Minister Benjamin Netanyahu ' s address to Congress for a convention raised questions about her stance on Israel amidst growing political tensions.
US \textbf{Vice} President Kamala Harris looks on as she visits the Reading Terminal Market in Philadelphia, Pennsylvania, US, July 13, 2024. ( photo credit : REUTERS/KEVIN MOHATT ) Custom dictates that when foreign leaders address a joint session of Congress, the Speaker of the House and the \textbf{vice} president sit behind them on the rostrum.
That was the case when Prime Minister Benjamin Netanyahu addressed Congress in 1996, and then-vice president Al Gore sat in one of the chairs behind him.
It was also true in 2011 when then-vice president Joe Biden took that chair.
Biden did not attend Netanyahu ' s speech in 2015, a clear sign of the White House ' s displeasure that this speech had been arranged despite their opposition.
On Wednesday, \textbf{Vice} President Kamala Harris was not in that seat during Netanyahu ' s address.
She should have been. ' Outragous and inexcusable ' @ @ @ @ @ @ @ @ @ @ attend the annual convention of a sorority being held in Indiana, is a bad look and a problematic first signal on Israel in her position as presumptive \textbf{presidential} \textbf{nominee}.
As House Speaker Mike Johnson put it,"It ' s outrageous to me and inexcusable that Kamala Harris is boycotting this joint session."Johnson, a \textbf{Republican}, was the lawmaker who invited Netanyahu, in part to accentuate the differences between Republicans and Democrats on Israel in an \textbf{election} year.
He knew many \textbf{Democrats} would boycott the speech.
Yet, the \textbf{Democratic} leadership signed off on the \textbf{invitation}, and as such, Harris should have shown up.
This would have been an opportunity to turn the tables on Johnson and demonstrate that there is no difference between the two \textbf{parties} when it comes to fundamental support for the Jewish state.
Harris must know that everything she says and does now will be carefully scrutinized.
Although her office said that she was absent from Netanyahu ' s speech because of a simple scheduling conflict, her absence @ @ @ @ @ @ @ @ @ @ speech by \textbf{Democratic} lawmakers.
Had Harris wanted to, she could have rearranged her schedule.
That she didn't, sends the signal that she was looking for an excuse -- as many other lawmakers did -- to not listen to Netanyahu.
And why would she take such a step?
To pander to the progressives of her \textbf{party} upset with the administration ' s overall strong support for Israel during the current war.
She needs to be savvy enough politically, however, to realize that as she is pandering to one flank of her \textbf{party}, she risks alienating pro-Israel \textbf{voters} on the other side.
Prefers a one-on-one Some, in Harris ' defense, will say that she is scheduled to meet Netanyahu for a private meeting during the prime minister ' s current trip, which should dispel any suggestions that she harbors him any enmity.
That meeting, however, is a private one in which the two leaders will learn where each other stands on various issues.
The speech in Congress symbolized the strong @ @ @ @ @ @ @ @ @ @ with the absence of other lawmakers -- sent a signal that there are cracks in those bonds, signals that Israel ' s enemies will detect and seek to exploit.
Another argument in her defense is that GOP \textbf{vice} \textbf{presidential} \textbf{candidate} J.D.
Vance also missed the speech, citing a previous \textbf{election} campaign commitment.
While true, there is a difference.
Vance has made his position crystal clear regarding the war in Gaza, Hamas ' s responsibility for the humanitarian disaster there, and the need for Israel to be allowed to defeat Hamas convincingly.
Harris has not.
On the contrary, since the beginning of the war, she has -- as even Israel ' s ambassador to the US Michael Herzog pointed out this week -- made some problematic comments regarding Israel.
Herzog said that overall, however, her record on Israel is positive.
Those"problematic"comments -- including accusations that Gazans were starving because of Israel and that anti-Israel protesters on campus were"showing exactly what the human emotion should be"-- raise @ @ @ @ @ @ @ @ @ @ : with Israel or the anti-Israel elements within the progressive wing of her \textbf{party}.
By attending Netanyahu ' s speech in the \textbf{vice} president ' s role as president of the senate, Harris would have sent one message.
By not attending, she sent another.

% matched lemmas: candidate, democrat, democratic, election, invitation, nominee, party, presidential, republican, vice, voter
\end{textsample}
